\subsection{In practice: twists in ROS 2, planners and pose graphs}\label{sec:pl:in-practice}

\subsubsection{The messages are twists}
A ROS~2 \texttt{geometry\_msgs/Twist} holds two 3-vectors, \texttt{linear} and \texttt{angular}. For a planar robot it is exactly this section's velocity twist, linear part first: $S=(\texttt{linear.x},\,\texttt{linear.y},\,\texttt{angular.z})$ per second, in the body frame. Velocity commands on \texttt{/cmd\_vel} are such twists. A differential-drive controller executes only \texttt{linear.x} and \texttt{angular.z}: the $v_y=0$ constraint of \S\ref{sec:pl:turn-centre}, which puts the turn centre on the axle. \texttt{nav\_msgs/Odometry} carries both a pose and a twist, and its definition states the frames precisely: ``the pose \dots\ in the coordinate frame given by \texttt{header.frame\_id}'' (usually \texttt{odom}), ``the twist \dots\ in the coordinate frame given by the \texttt{child\_frame\_id}'' (usually \texttt{base\_link}). The pose is $T_{\text{odom},\text{base}}$, and the twist is a \emph{body} twist. The controller that produces it integrates wheel increments with the exact arc of \S\ref{sec:pl:twists} and also publishes \texttt{odom}$\to$\texttt{base\_link} on tf (\S\ref{sec:rot:tf-map-odom}).

\subsubsection{Planners are built from exponentials}
The shortest path for a car that drives forward with a minimum turning radius $r$ (Dubins, 1957) is a sequence of at most three segments, each a left arc, a straight line or a right arc. Allowing reverse as well (Reeds and Shepp, 1990) gives at most five. Each segment is the exponential of a constant unit twist:
\[
L(a)=\exp\big(a[(r,0,1)]\big),\qquad S(d)=\exp\big(d[(1,0,0)]\big),\qquad R(a)=\exp\big(a[(r,0,-1)]\big),
\]
so evaluating a candidate path is a product of three closed-form matrices.

\begin{example}{a Dubins path in three lines}
With $r=0.4$\,m (Nav2's default \texttt{minimum\_turning\_radius}), the path ``left $90^\circ$, straight 1\,m, left $90^\circ$'' ends at
\[
L(\tfrac\pi2)\,S(1)\,L(\tfrac\pi2):\quad(0,\ 1.8)\,\text{m},\ \text{heading }180^\circ,\quad\text{length }2.257\,\text{m};
\]
``left, straight, right'' ends at $(0.8,1.8)$\,m, heading $0^\circ$. Sanity check by hand: the first arc ends at $(r,r)$ facing up, the straight adds 1\,m in $y$ to reach $(0.4,1.4)$, and the last quarter arc turns about $(0,1.4)$ to end at $(0,1.8)$ facing $-x$. Nav2's Smac Hybrid-A* planner searches over exactly these primitives: its \texttt{motion\_model\_for\_search} is \texttt{DUBIN} or \texttt{REEDS\_SHEPP}, and successors of a search node are generated by short arcs and straights, i.e.\ by exponentials of a handful of twists.
\end{example}

\subsubsection{Estimators use the logarithm}
Graph SLAM back-ends for planar robots (Karto and \texttt{slam\_toolbox} in ROS~2) store relative-pose measurements between scans and minimise residuals of the form $\vexa(\log(Z_{ab}^{-1}T_a^{-1}T_b))$ from \S\ref{sec:pl:logarithms}. The Lie-group libraries used to build such systems (GTSAM's \texttt{Pose2::Expmap}/\texttt{Logmap}, Sophus' \texttt{SE2}, manif's \texttt{SE2}) implement precisely the closed forms derived here, the $V(\theta)$ factor included, and all three order the planar tangent vector as $(x,y,\theta)$, linear part first. Solà, Deray and Atchuthan's \emph{A micro Lie theory for state estimation in robotics} (2018) is the standard short reference that connects this section to those libraries.
